OCC-RAG models -- our SLM family specifically designed for Context QA should possess the following properties: (1) multi‑hop inference and commonsense reasoning for complex questions; (2) avoidance of memorization (faithfulness to context, no conflict with internal knowledge); and (3) safe abstention when the provided context is insufficient.

Mid-training serves as a core stage that explicitly shapes the SLM's reasoning architecture for Context QA, enabling developers to gain fine-grained control over evidence combination and yielding more reliable, interpretable, and context-aligned downstream QA. Large-scale agentic mid-training on synthetic trajectories across math, code, and tool-use further internalizes planning and reflection, unlocking native agentic potential in lightweight models and surpassing larger baselines on agentic benchmarks~\citep{pleias,youtu}.

\paragraph{Mid‑training enables strong multi‑hop reasoning} 
Mid‑training on reasoning‑trace datasets improves QA performance of SLMs by training them on the functional structure of multi‑hop inference -- such as subquestion decomposition, information retrieval, and step‑wise verification -- rather than simply teaching them to reproduce surface‑level answer patterns. This ``structural'' signal helps SLMs internalize the process by which correct answers are reached, which in turn boosts generalization to new QA instances and reduces reliance on superficial shortcuts~\citep{lee-hockenmaier-2025-evaluating,liang2026longdocument}.

\textbf{Mid-training supports faithful, context-grounded, non-memorized QA} Mid-training on context-based reasoning traces that always tie each step back to the provided text—ensuring strict faithfulness to evidence helps SLMs learn to solve QA tasks without memorizing facts or hallucinating. Studies on multi-hop QA show that fine-tuning alone on raw text or continual pretraining yields only limited gains, whereas structured or supervised mid-training on evidence-anchored traces substantially improves answer accuracy without relying on internal knowledge~\citep{ren2026sinbenchtracingnativeevidence,li-etal-2024-teaching}.

\textbf{Mid‑training encourages calibrated abstention} When mid‑training includes ``context‑insufficient'' or unanswerable examples annotated with explicit reasoning‑trace patterns (e.g., ``no evidence in context''), the SLM learns to recognize when the context does not support a confident answer. This kind of structured‑reasoning training is shown to improve the model’s ability to abstain appropriately, rather than hallucinate, on information‑limited tasks and partial‑context environments. In effect, mid‑training turns abstention into a learned reasoning behavior rather than a heuristic, making the SLM more reliable in high‑stakes QA settings~\citep{zhou2026when,wen-etal-2024-characterizing}.